\documentclass{report}
\usepackage{amsmath,amssymb}
\setlength{\parindent}{0mm}
\setlength{\parskip}{1em}
\begin{document}
\begin{center}
\rule{6in}{1pt} \
{\large Xin Wang \\
{\bf Lanczos-filter subspace iteration for self-consistent-field calculation}}

311 Stevens Circle \\ Apt3A \\ Aberdeen \\ MD 21001
\\
{\tt xin.c.wang2.ctr@mail.mil}\\
Xin C Wang\\
Kenneth Leiter\\
Alexander Breuer\\
Jaroslaw Knap\end{center}

Kohn-Sham density functional theory (K-S DFT) is widely-used for
ground-state electronic structure calculations. Its bottleneck is the
computation of N occupied eigen-states of a non-linear eigenvalue problem, where
N is the number of the electron-pairs. The non-linear eigenvalue problem is
solved through a self-consistent-field (SCF) iteration where at each iteration
a linear eigenvalue problem is solved.

Due to the large number of electrons in
the model problem and the high accuracy (around $1e-4$ relative error) required
for K-S DFT, the resulting degrees of freedom in the matrices can range from
hundreds of thousands to tens of millions. Furthermore, the condition number
of the matrices can be poor depending on the choice of the discretization
basis. For real-space
based methods such as finite difference or finite elements, the condition number
can be in the range of millions for DFT problems that
model all the electrons. These characteristics of the DFT problem make it very
difficult for conventional iterative eigensolvers to find the wanted eigen-states.

Zhou \emph{et al.} (Journal of Computational Physics,
2006) proposed the Chebyshev-filtered subspace iteration (CheFSI) as a
non-linear subspace iteration method to extract the wanted eigen-subspace.
CheFSI starts with an initial guess to the wanted eigen-subspace of the initial
matrix $H_{\rm initial}$, and progressively through the SCF iterations
transform it to the wanted eigen-subspace of the final converged matrix
$H_{\rm final}$ using a
Chebyshev-polynomial filter. We propose a Lanczos-filter subspace
iteration (LanFSI) that uses an interpolative-polynomial approximation to
smeared step function from the Lanczos iteration to filter the wanted
eigen-subspace. We observe that for matrices with
a large condition number, LanFSI is a more efficient filter than CheFSI.
Consequently, LanFSI achieves convergence with fewer SCF iterations than
CheFSI using the same polynomial degree in all-electron calculations.


\end{document}
